\subsection{The Springer correspondence}\label{ss:Spr corr} Let $W$ be the Weyl group of $G$. In 1976, Springer \cite{Springer} made the fundamental observation that there is natural $W$-action on $\cohog{*}{\sB_{e}}$, even though $W$ does not act on $\sB_{e}$ as automorphisms of varieties. 


\begin{theorem}[Springer {\cite[Thm 6.10]{Springer}}]\label{th:Spr corr} 
\begin{enumerate}
\item For each nilpotent element $e$, there is a natural graded action of $W$ on $\cohog{*}{\sB_{e}}$ that commutes with the action of $A_{e}$.
\item For each nilpotent element $e$ and each irreducible representation $\rho$ of $A_{e}$, the multiplicity space $M(e,\rho):=\Hom_{A_{e}}(\rho, \cohog{2d_{e}}{\sB_{e}})$ is either zero or an irreducible representation of $W$ under the action in part (1).
\item Each irreducible representation $\chi$ of $W$ appears as $M(e,\rho)$ for a unique pair $(e,\rho)$ up to $G$-conjugacy. The assignment $\chi\mapsto (e,\rho)$ thus gives an injection
\begin{equation}\label{Spr corr}
\Irr(W)\incl \{(e,\rho)\}/G.
\end{equation}
\end{enumerate}
\end{theorem}

\subsubsection{Convention}\label{sss:two actions} In fact there are two natural actions of $W$ on $\cohog{*}{\sB_{e}}$ that differ by tensoring with the sign representation of $W$. In these notes we use the action that is normalized by the following properties. The trivial representation of $W$ corresponds to regular nilpotent $e$ and the trivial $\rho$. The sign representation of $W$ corresponds to $e=0$. Note however that Springer's original paper \cite{Springer} uses the other action.

\subsubsection{The case $e=0$} Taking $e=0$, Springer's theorem gives a graded action of $W$ on $\cohog{*}{\calB}$. What is this action? First, this action can be seen geometrically by considering $G/T$ instead of $\sB=G/B$. In fact, since $N_{G}(T)/T=W$, the right action of $N_{G}(T)$ on $G/T$ induces an action of $W$ on $G/T$, which then induces an action of $W$ on $\cohog{*}{G/T}$. Since the projection $G/T\to G/B$ is an affine space bundle, it follows that $\cohog{*}{\sB}\cong\cohog{*}{G/T}$. It can be shown that under this isomorphism, the action of $W$ on $\cohog{*}{G/T}$ corresponds exactly to Springer's action on $\cohog{*}{\calB}$.

Let $S=\Sym(\xch(T)\otimes\Ql)$ be the graded symmetric algebra where $\xch(T)$ has degree $2$. The reflection representation of $W$ on $\xch(T)$ then induces a graded action of $W$ on $S$. Recall Borel's presentation of the cohomology ring of the flag variety
\begin{equation}\label{Borel HB}
\cohog{*}{\calB,\Ql}\cong S/(S^{W}_{+})
\end{equation}
where $S_{+}\subset S$ is the ideal spanned by elements of positive degree, and $(S^{W}_{+})$ denotes the ideal of $S$ generated by $W$-invariants on $S_{+}$. Then \eqref{Borel HB} is in fact an isomorphism of $W$-modules (see \cite[Prop 7.2]{Springer}). By a theorem of Chevalley, $S/(S^{W}_{+})$ is isomorphic to the regular representation of $W$, therefore, as a $W$-module, $\cohog{*}{\sB}$ is also isomorphic to the regular representation of $W$.  

\begin{remark} The target set in \eqref{Spr corr} can be canonically identified with the set of isomorphism classes of irreducible $G$-equivariant local systems on nilpotent orbits. In fact, for an irreducible $G$-equivariant local systems $\calL$ on a nilpotent orbit $\calO\subset\sN$, its stalk at $e\in\calO$ gives an irreducible representation $\rho$ of the centralizer $C_{G}(e)$ which factors through $A_{e}$. Note that the notion of $G$-equivariance changes when $G$ varies in a fixed isogeny class.  It is possible to extend the above injection \eqref{Spr corr} into a bijection by supplementing $\Irr(W)$ with $\Irr(W')$ for a collection of smaller Weyl groups. This is called the {\em generalized Springer correspondence} discovered by Lusztig \cite{LuIC}.
\end{remark}

Springer's original proof of Theorem \ref{th:Spr corr} uses trigonometric sums over $\frg(\FF_{q})$ and, when $k$ has characteristic zero, his proof uses reduction to finite fields. The following theorem due to Borho and MacPherson \cite{BM} can be used to give a direct proof of the Springer correspondence for all base fields $k$ of large characteristics or characteristic zero. To state it, we need to use the language of constructible (complexes of) $\Ql$-sheaves and perverse sheaves, for which we refer to the standard reference \cite{BBD} and de Cataldo's lectures \cite{dC} in this volume. 

\begin{theorem}\label{th:BM} The complex $\calS:=\bR\pi_{*}\Ql[\dim\sN]$ is a perverse sheaf on $\sN$ whose endomorphism ring is canonically isomorphic to the group algebra $\Ql[W]$. In particular, $W$ acts on the stalks of $\bR\pi_{*}\Ql$, i.e., $W$ acts on $\cohog{*}{\sB_{e}}$ for all $e\in\sN$.
\end{theorem}


We sketch three constructions of the $W$-action on $\bR\pi_{*}\Ql[\dim\sN]$.

\subsubsection{Construction via middle extension}\label{sss:mid} This construction (or rather the version where $\frg$ is replaced by $G$) is due to Lusztig \cite[\S3]{LuGreen}. The dimension formula for Springer fibers \eqref{Spr dim} imply that 
\begin{itemize}
\item The map $\pi:\tN\to\sN$ is semismall. \footnote{A proper surjective map $f:X\to Y$ of irreducible varieties is called {\em semismall} (resp. {\em small}) if for any $d\geq 1$, $\{y\in Y|\dim f^{-1}(y)\geq d\}$ has codimension at least $2d$ (resp. $2d+1$) in $Y$.}
\end{itemize}
There is an extension of the dimension formula \eqref{Spr dim} for the dimension of $\tB_{X}$ valid for all elements $X\in \frg$. Using this formula one can show that
\begin{itemize}
\item The map $\pi_{\frg}:\tg\to\frg$ is small.
\end{itemize}
As a well-known fact in the theory of perverse sheaves, the smallness of $\pi_{\frg}$ (together with the fact that $\tg$ is smooth) implies that $\calS_{\frg}:=\bR\pi_{\frg,*}\Ql[\dim\frg]$ is a perverse sheaf which is  the middle extension of its restriction to any open dense subset of $\frg$. Over the regular semisimple locus $\frg^{\rs}\subset \frg$, $\pi_{\frg}$ is a $W$-torsor, therefore $\calS_{\frg}|_{\frg^{rs}}$ is a local system shifted in degree $-\dim\frg$ that admits an action of $W$. By the functoriality of middle extension, $\calS_{\frg}$ admits an action of $W$. Taking stalks of $\calS_{\frg}$, we get an action of $W$ on $\cohog{*}{\tB_{X}}$ for all $X\in\frg$. 

In particular, for a nilpotent element $e$, we get an action of $W$ on $\cohog{*}{\sB_{e}}=\cohog{*}{\tB_{e}}$, because $\tB_{e}$ and $\sB_{e}$ have the same reduced structure. This is the action defined by Springer in his original paper \cite{Springer}, {\em which differs from our action by tensoring with the sign character of $W$}.

